\begin{table}[!t]
\centering
\caption{Principales notations du modèle MCAD.}
\scriptsize
\begin{tabularx}{\columnwidth}{@{}p{2.35cm}X@{}}
\toprule
Symbole & Description \\
\midrule
$F$ & Ensemble des tables de faits \\
$D$ & Ensemble des dimensions \\
$L(D)$ & Niveaux de la dimension $D$ \\
$g$ & Grain d'analyse (produit de niveaux) \\
$M$ & Ensemble des mesures et KPI \\
CKG & Graphe de connaissances contextuel \\
$O$ & Objectif stratégique actif \\
$C^*(O)$ & Ensemble des contraintes de $O$ \\
$N(c)$ & Nœuds virtuels associés à la contrainte $c$ \\
$QP$ & Plan de requête canonique \\
$\mathrm{Real}(QP)$ & Nœuds virtuels réalisables par $QP$ \\
$C_{\mathrm{eval}}(QP,O)$ & Contraintes complètement calculables via $QP$ \\
$E_t$ & Évidence en nœuds virtuels acquise et retenue à l'étape $t$ \\
$C_Q^{\le t}(O)$ & Accumulation par requêtes des contraintes complétées individuellement \\
$C_E^{\le t}(O)$ & Contraintes calculables à partir de l'évidence cumulée $E_t$ \\
$\kappa_c(X)$ & Degré de calculabilité intra-contrainte de $c$ \\
$\psi(QP,O)$ & Calculabilité graduée induite par $QP$ \\
$\varphi(QP,O)$ & Couverture des contraintes complètement calculables \\
\bottomrule
\end{tabularx}
\label{tab:notation}
\normalsize
\end{table}
